% problem_id: S001-lemma-serre-subcategory-K-groups
% source_id: S001 (Stacks Project)
% source_file: homology.tex
% statement_locator: homology.tex lines 2400--2425
% proof_locator: homology.tex lines 2427--2614
% env: lemma
% id_note: label 在本来源内唯一
% pairing: tight (verified)
% status: complete (statement + author proof verbatim + reason + locators)
%
\begin{lemma}
\label{lemma-serre-subcategory-K-groups}
Let $\mathcal{A}$ be an abelian category.
Let $\mathcal{C} \subset \mathcal{A}$ be a Serre subcategory and
set $\mathcal{B} = \mathcal{A}/\mathcal{C}$.
\begin{enumerate}
\item The exact functors $\mathcal{C} \to \mathcal{A}$ and
$\mathcal{A} \to \mathcal{B}$ induce an exact sequence
$$
K_0(\mathcal{C}) \to
K_0(\mathcal{A}) \to
K_0(\mathcal{B}) \to
0
$$
of $K$-groups, and
\item the kernel of $K_0(\mathcal{C}) \to K_0(\mathcal{A})$ is equal
to the collection of elements of the form
$$
[H^0(M, \varphi, \psi)] - [H^1(M, \varphi, \psi)]
$$
where $(M, \varphi, \psi)$ is a complex as in (\ref{equation-cyclic-complex})
with the property that it becomes exact in $\mathcal{B}$; in other words
that $H^0(M, \varphi, \psi)$ and $H^1(M, \varphi, \psi)$ are
objects of $\mathcal{C}$.
\end{enumerate}
\end{lemma}

\begin{proof}
Proof of (1). It is clear that $K_0(\mathcal{A}) \to K_0(\mathcal{B})$
is surjective and that the composition $K_0(\mathcal{C}) \to
K_0(\mathcal{A}) \to K_0(\mathcal{B})$ is zero. Let $x \in K_0(\mathcal{A})$
be an element mapping to zero in $K_0(\mathcal{B})$. We can write
$x = [A] - [A']$ with $A, A'$ in $\mathcal{A}$ (fun exercise).
Denote $B, B'$ the corresponding objects of $\mathcal{B}$. The fact that
$x$ maps to zero in $K_0(\mathcal{B})$
means that there exists a finite set $I = I^+ \amalg I^{-}$,
for each $i \in I$ a short exact sequence
$$
0 \to B_i \to B'_i \to B''_i \to 0
$$
in $\mathcal{B}$ such that we have
$$
[B] - [B'] = \sum\nolimits_{i \in I^{+}} ([B'_i] - [B_i] - [B''_i])
-
\sum\nolimits_{i \in I^{-}} ([B'_i] - [B_i] - [B''_i])
$$
in the free abelian group on isomorphism classes of objects of $\mathcal{B}$.
We can rewrite this as
$$
[B]
+ \sum\nolimits_{i \in I^{+}} ([B_i] + [B''_i])
+ \sum\nolimits_{i \in I^{-}} [B'_i]
=
[B']
+ \sum\nolimits_{i \in I^{-}} ([B_i] + [B''_i])
+ \sum\nolimits_{i \in I^{+}} [B'_i].
$$
Since the right and left hand side should contain the same isomorphism classes
of objects of $\mathcal{B}$ counted with multiplicity, this means there should
be a bijection
$$
\tau :
\{B\} \amalg \{B_i, B''_i; i \in I^+\} \amalg \{B'_i; i \in I^-\}
\longrightarrow
\{B'\} \amalg \{B_i, B''_i; i \in I^-\} \amalg \{B'_i; i \in I^+\}
$$
such that $N$ and $\tau(N)$ are isomorphic in $\mathcal{B}$.
The proof of Lemmas \ref{lemma-serre-subcategory-is-kernel} and
\ref{lemma-localization-abelian} show that we choose for $i \in I$
a short exact sequence
$$
0 \to A_i \to A'_i \to A''_i \to 0
$$
in $\mathcal{A}$ such that $B_i, B'_i, B''_i$ are isomorphic to the
images of $A_i, A'_i, A''_i$ in $\mathcal{B}$. This implies that the
corresponding bijection
$$
\tau :
\{A\} \amalg \{A_i, A''_i; i \in I^+\} \amalg \{A'_i; i \in I^-\}
\longrightarrow
\{A'\} \amalg \{A_i, A''_i; i \in I^-\} \amalg \{A'_i; i \in I^+\}
$$
satisfies the property that $M$ and $\tau(M)$ are objects of $\mathcal{A}$
which become isomorphic in $\mathcal{B}$. This means $[M] - [\tau(M)]$
is in the image of $K_0(\mathcal{C}) \to K_0(\mathcal{A})$. Namely,
the isomorphism in $\mathcal{B}$ is given by a diagram
$M \leftarrow M' \rightarrow \tau(M)$ in $\mathcal{A}$ where both
$M' \to M$ and $M' \to \tau(M)$ have kernel and cokernel in $\mathcal{C}$.
Working backwards we conclude that $x = [A] - [A']$ is in the image
of $K_0(\mathcal{C}) \to K_0(\mathcal{A})$ and the proof of part (1)
is complete.

\medskip\noindent
Proof of (2). The proof is similar to the proof of (1) but slightly
more bookkeeping is involved. First we remark that any class of the type
$[H^0(M, \varphi, \psi)] - [H^1(M, \varphi, \psi)]$ is zero
in $K_0(\mathcal{A})$ by the following calculation
\begin{align*}
0 & = [M] - [M] \\
& =  [\Ker(\varphi)] + [\Im(\varphi)]
- [\Ker(\psi)] - [\Im(\psi)] \\
& =
[\Ker(\varphi)/\Im(\psi)] -
[\Ker(\psi)/\Im(\varphi)] \\
& = [H^1(M, \varphi, \psi)] - [H^0(M, \varphi, \psi)]
\end{align*}
as desired. Hence it suffices to show that any element in the kernel
of $K_0(\mathcal{C}) \to K_0(\mathcal{A})$ is of this form.

\medskip\noindent
Any element $x$ in $K_0(\mathcal{C})$ can be represented as the
difference $x = [P] - [Q]$ of two objects of $\mathcal{C}$ (fun exercise).
Suppose that this element maps to zero in $K_0(\mathcal{A})$.
This means that there exist
\begin{enumerate}
\item a finite set $I = I^{+} \amalg I^{-}$,
\item for $i \in I$ a short exact sequence $0 \to A_i \to B_i \to C_i \to 0$
in $\mathcal{A}$
\end{enumerate}
such that
$$
[P] - [Q] =
\sum\nolimits_{i \in I^{+}} ([B_i] - [A_i] - [C_i])
-
\sum\nolimits_{i \in I^{-}} ([B_i] - [A_i] - [C_i])
$$
in the free abelian group on the objects of $\mathcal{A}$.
We can rewrite this as
$$
[P]
+ \sum\nolimits_{i \in I^{+}} ([A_i] + [C_i])
+ \sum\nolimits_{i \in I^{-}} [B_i]
=
[Q]
+ \sum\nolimits_{i \in I^{-}} ([A_i] + [C_i])
+ \sum\nolimits_{i \in I^{+}} [B_i].
$$
Since the right and left hand side should contain the same objects
of $\mathcal{A}$ counted with multiplicity, this means there should be
a bijection $\tau$ between the terms which occur above. Set
$$
T^{+} =
\{p\}\ \amalg\ \{a, c\} \times I^{+}\ \amalg\ \{b\} \times I^{-}
$$
and
$$
T^{-} =
\{q\}\ \amalg\ \{a, c\} \times I^{-}\ \amalg\ \{b\} \times I^{+}.
$$
Set $T = T^{+} \amalg T^{-} = \{p, q\} \amalg \{a, b, c\} \times I$.
For $t \in T$ define
$$
O(t)
=
\left\{
\begin{matrix}
P & \text{if} & t = p \\
Q & \text{if} & t = q \\
A_i & \text{if} & t = (a, i) \\
B_i & \text{if} & t = (b, i) \\
C_i & \text{if} & t = (c, i)
\end{matrix}
\right.
$$
Hence we can view $\tau : T^{+} \to T^{-}$ as a bijection
such that $O(t) = O(\tau(t))$ for all $t \in T^{+}$.
Let $t^{-}_0 = \tau(p)$ and let $t^{+}_0 \in T^{+}$ be the
unique element such that $\tau(t^{+}_0) = q$.
Consider the object
$$
M^{+} = \bigoplus\nolimits_{t \in T^{+}} O(t)
$$
By using $\tau$ we see that it is equal to the object
$$
M^{-} = \bigoplus\nolimits_{t \in T^{-}} O(t)
$$
Consider the map
$$
\varphi : M^{+} \longrightarrow M^{-}
$$
which on the summand $O(t) = A_i$ corresponding to $t = (a, i)$, $i \in I^{+}$
uses the map $A_i \to B_i$ into the summand $O((b, i)) = B_i$ of $M^{-}$
and on the summand $O(t) = B_i$ corresponding to $(b, i)$, $i \in I^{-}$
uses the map $B_i \to C_i$ into the summand $O((c, i)) = C_i$ of $M^{-}$.
The map is zero on the summands corresponding to $p$
and $(c, i)$, $i \in I^{+}$.
Similarly, consider the map
$$
\psi : M^{-} \longrightarrow M^{+}
$$
which on the summand $O(t) = A_i$ corresponding to $t = (a, i)$, $i \in I^{-}$
uses the map $A_i \to B_i$ into the summand $O((b, i)) = B_i$ of $M^{+}$
and on the summand $O(t) = B_i$ corresponding to $(b, i)$, $i \in I^{+}$
uses the map $B_i \to C_i$ into the summand $O((c, i)) = C_i$ of $M^{+}$.
The map is zero on the summands corresponding to $q$ and
$(c, i)$, $i \in I^{-}$.

\medskip\noindent
Note that the kernel of $\varphi$ is equal to the direct sum of the
summand $P$ and the summands $O((c, i)) = C_i$, $i \in I^{+}$ and
the subobjects $A_i$ inside the summands $O((b, i)) = B_i$, $i \in I^{-}$.
The image of $\psi$ is equal to the direct sum of the
summands $O((c, i)) = C_i$, $i \in I^{+}$ and
the subobjects $A_i$ inside the summands $O((b, i)) = B_i$, $i \in I^{-}$.
In other words we see that
$$
P \cong \Ker(\varphi)/\Im(\psi).
$$
In exactly the same way we see that
$$
Q \cong \Ker(\psi)/\Im(\varphi).
$$
Since as we remarked above the existence of the bijection
$\tau$ shows that $M^{+} = M^{-}$ we see that the lemma follows.
\end{proof}
